% !TeX program = xelatex
% Απαιτεί XeLaTeX (λόγω fontspec/polyglossia) -- ΔΕΝ μεταγλωττίζεται με pdflatex.
\documentclass[11pt,a4paper]{article}

\usepackage{fontspec}
\usepackage{polyglossia}
\setmainlanguage{greek}
\setotherlanguage{english}
\setmainfont{Linux Libertine O}
\setsansfont{Linux Biolinum O}
\setmonofont{DejaVu Sans Mono}[Scale=0.9]

\usepackage[a4paper,margin=2.6cm]{geometry}
\usepackage{parskip}
\usepackage{amsmath,amssymb}
\usepackage{booktabs}
\usepackage{array}
\usepackage{graphicx}
\graphicspath{{../results/}}
\usepackage{enumitem}
\usepackage{caption}
\usepackage[hidelinks,colorlinks=true,linkcolor=blue!50!black,citecolor=blue!50!black,urlcolor=blue!50!black]{hyperref}
\usepackage{xcolor}

\newcommand{\file}[1]{\texttt{#1}}

\let\oldunderscore\_
\renewcommand{\_}{\oldunderscore\allowbreak}
\setlength{\emergencystretch}{3em}

\setlist[itemize]{topsep=2pt,itemsep=2pt}
\setlist[enumerate]{topsep=2pt,itemsep=2pt}

\title{\textbf{Online Convex Optimization based tuning of\\ Mixture of Experts for Sports Outcome Prediction}\\[6pt]
\large Σύντομη αναφορά προόδου}
\author{[Ονοματεπώνυμο φοιτητή]}
\date{\today}

\begin{document}
\maketitle

% =====================================================================
\section{Τι κάνει η εργασία}
% =====================================================================

Κάθε στοιχηματική εταιρεία δίνει τις δικές της πιθανότητες $(p_H, p_D, p_A)$ για το
αποτέλεσμα ενός ποδοσφαιρικού αγώνα (Νίκη Γηπεδούχου / Ισοπαλία / Νίκη Φιλοξενούμενου).
Αντιμετωπίζουμε κάθε bookmaker ως έναν \emph{expert}, και μαθαίνουμε \emph{online}
(αγώνα προς αγώνα, χωρίς να ξαναδούμε το παρελθόν) πώς να τους συνδυάζουμε σε ένα ενιαίο
μείγμα που να τους ξεπερνά όλους μαζί. Χρησιμοποιούμε τρεις κλασικούς αλγορίθμους Online
Convex Optimization (Hedge, OGD, FTRL), γενικευμένους ώστε να δουλεύουν και όταν κάποιοι
bookmakers δεν δίνουν απόδοση σε κάθε αγώνα (μερική κάλυψη).

% =====================================================================
\section{Δεδομένα}
% =====================================================================

Πηγή: football-data.co.uk, ενότητα "Main Leagues". \textbf{22 διοργανώσεις}
(Αγγλία $\times$5, Σκωτία $\times$4, Γερμανία/Ιταλία/Ισπανία/Γαλλία $\times$2,
Ολλανδία/Βέλγιο/Πορτογαλία/Τουρκία/Ελλάδα $\times$1), \textbf{10 σεζόν} (2016/17--2025/26),
\textbf{76.639 αγώνες}, εκ των οποίων 76.584 έχουν τουλάχιστον μία χρήσιμη απόδοση.
Μετά την αφαίρεση μη-ανεξάρτητων/aggregate στηλών (π.χ.\ Max/Avg) και τη συγχώνευση ενός
γνωστού rebrand (VC Bet $\to$ BetVictor), μένουν \textbf{26 ανεξάρτητοι bookmakers}.

% =====================================================================
\section{Μεθοδολογία, σε συντομία}
% =====================================================================

\begin{itemize}
  \item \textbf{De-vig}: οι raw αποδόσεις μετατρέπονται σε καθαρές πιθανότητες
        (proportional normalization) αφαιρώντας το περιθώριο κέρδους (overround) του
        κάθε bookmaker.
  \item \textbf{Sleeping Experts / Specialists reduction}: επιτρέπει στο Hedge/OGD/FTRL να
        δουλεύουν σωστά ακόμη κι όταν ένας bookmaker "κοιμάται" (δεν δίνει απόδοση) σε
        κάποιους αγώνες -- επανακανονικοποίηση μόνο πάνω στους ενεργούς κάθε γύρο, με το
        state των κοιμισμένων παγωμένο.
  \item \textbf{Βαθμονόμηση (calibration)}: δεύτερο στάδιο OCO (online temperature scaling)
        πάνω στο μείγμα, για διόρθωση του γνωστού \emph{favorite-longshot bias} (οι αγορές
        στοιχημάτων τείνουν να υπερτιμούν τα "μεγάλα" σενάρια και να υποτιμούν τα φαβορί).
  \item \textbf{Στατιστική σημαντικότητα}: moving block bootstrap (όχι απλό t-test) --
        λαμβάνει υπόψη ότι οι διαδοχικές ζημίες ανά αγώνα δεν είναι ανεξάρτητες.
  \item \textbf{Οικονομικός έλεγχος}: προσομοίωση στοιχηματισμού Kelly, leave-one-
        bookmaker-out (ο bookmaker εναντίον του οποίου στοιχηματίζουμε αποκλείεται πρώτα
        από το μείγμα, ώστε να μην είναι κυκλικό).
\end{itemize}

Καμία θεωρητική απόδειξη δεν παρατίθεται εδώ -- υπάρχουν αναλυτικά στο πλήρες κείμενο
(\file{thesis\_report.tex}).

% =====================================================================
\section{Κύρια ευρήματα}
% =====================================================================

\subsection{Πυρήνας: Hedge / OGD / FTRL έναντι bookmakers}

\begin{table}[h]
\centering
\begin{tabular}{@{}lrr@{}}
\toprule
\textbf{Μέθοδος} & \textbf{Raw log-loss} & \textbf{Βαθμονομημένο log-loss} \\
\midrule
PSC (Pinnacle, closing) & 0,9971 & \textbf{0,9968} \\
OGD (καλύτερος αλγόριθμος) & 0,9983 & 0,9976 \\
Hedge & 0,9985 & 0,9978 \\
FTRL & 0,9985 & 0,9978 \\
Uniform average & 0,9989 & 0,9981 \\
\bottomrule
\end{tabular}
\caption{Tier-A κατάταξη (κάλυψη $\geq$70\%), ταξινομημένη κατά βαθμονομημένο log-loss.
Χαμηλότερο = καλύτερο.}
\end{table}

\begin{itemize}
  \item Το \textbf{OGD} είναι ο καλύτερος από τους τρεις αλγορίθμους, και μαζί με το Hedge
        και το FTRL ξεπερνούν σημαντικά το naive uniform average.
  \item Ο \textbf{PSC} (Pinnacle, closing odds) έχει στατιστικά σημαντικό προβάδισμα έναντι
        του OGD -- κανένας μεμονωμένος αλγόριθμος δεν ξεπερνά τον καλύτερο "sharp"
        bookmaker, μόνο τους υπόλοιπους.
  \item Η \textbf{βαθμονόμηση βοηθάει σημαντικά}, με συνέπεια, σε όλους τους αλγορίθμους
        και bookmakers.
  \item Οι \textbf{αποδόσεις κλεισίματος} (closing) δίνουν σημαντικά καλύτερο μείγμα από
        τις αποδόσεις ανοίγματος (opening) -- αναμενόμενο, αφού η αγορά έχει απορροφήσει
        περισσότερη πληροφορία μέχρι το κλείσιμο.
\end{itemize}

\subsection{Επεκτάσεις (μία γραμμή η καθεμία)}

\begin{itemize}
  \item \textbf{Fixed-Share}: διορθώνει bookmakers με τεχνητά διογκωμένο τελικό βάρος
        (π.χ.\ ενεργοί σε μία μόνο σεζόν), με μικρό κόστος στο συνολικό log-loss.
  \item \textbf{Contextual experts (ανά λίγκα)}: \emph{δεν} βοηθάει -- το ενιαίο μείγμα
        ξεπερνά τα specialists ανά λίγκα (πλην Σκωτίας, οριακά).
  \item \textbf{Value betting}: καμία σημαντική άκρη έναντι σκληρού bookmaker (B365C)·
        μεγάλη, σημαντική άκρη έναντι πιο αδύναμου (BW) -- με την επιφύλαξη ότι
        πραγματικά όρια λογαριασμού δεν προσομοιώνονται.
  \item \textbf{Ακρίβεια ανά κατηγορία}: οι ισοπαλίες είναι ουσιαστικά απρόβλεπτες ως
        \emph{argmax} πρόβλεψη ($<$0,7\% ακρίβεια), παρότι $\sim$26\% των αγώνων.
  \item \textbf{Χρονικές τάσεις (10 σεζόν)}: καμία σημαντική τάση στο log-loss/ECE, αλλά
        το overround \emph{αυξάνεται} σημαντικά ($\sim$4,4\%$\to$6,0\%).
  \item \textbf{Per-season βαθμονόμηση / warm-start}: σημαντικά καλύτερα από raw, αλλά
        στατιστικά αδιάκριτα από το ενιαίο global calibration.
\end{itemize}

% =====================================================================
\section{Πρόσφατη δουλειά: εναλλακτική μέθοδος de-vig (Shin)}
% =====================================================================

Ελέγξαμε αν τα παραπάνω ευρήματα είναι artifact της επιλογής μεθόδου de-vig
(proportional normalization), δοκιμάζοντας τη μέθοδο του \textbf{Shin (1992)} --
μοντελοποιεί μια αναλογία "insider" στοιχηματιστών και ανακατανέμει το περιθώριο του
bookmaker υπέρ του φαβορί, αντί να το μοιράζει ομοιόμορφα.

\begin{itemize}
  \item Η \textbf{κατάταξη δεν αλλάζει}: ο PSC παραμένει πρώτος, με την ίδια σειρά
        αλγορίθμων/bookmakers, είτε με proportional είτε με Shin de-vig.
  \item Στο \textbf{raw} στάδιο, το Shin είναι στατιστικά σημαντικά καλύτερο
        (log-loss \emph{και} Brier) σε όλες τις 11 οντότητες (4 αλγόριθμοι + 7 Tier-A
        bookmakers), $p<0{,}0001$ παντού.
  \item Η \textbf{ποιότητα βαθμονόμησης} (Expected Calibration Error) βελτιώνεται
        δραματικά με Shin -- μειώνεται κατά $\sim$35--57\% σε κάθε σειρά (π.χ.\ OGD:
        0,0091 $\to$ 0,0043).
  \item Μετά όμως τη βαθμονόμηση, το πλεονέκτημα του Shin \textbf{συρρικνώνεται ή
        εξαφανίζεται} -- και, αντίστροφα, το ίδιο το calibration παύει να είναι
        στατιστικά σημαντικό για τα περισσότερα series \emph{μέσα} στο Shin (ενώ είναι
        σημαντικό παντού μέσα στο proportional normalization). Ερμηνεία: οι δύο μέθοδοι
        διορθώνουν το \emph{ίδιο} favorite-longshot bias, με διαφορετικό μηχανισμό --
        μόλις γίνει η μία διόρθωση, η άλλη γίνεται σε μεγάλο βαθμό περιττή.
  \item Το ίδιο μοτίβο επιβεβαιώνεται και στην οικονομική προσομοίωση (value betting).
\end{itemize}

\begin{figure}[h]
\centering
\includegraphics[width=0.95\textwidth]{shin_vs_basic_calibration_reliability.png}
\caption{Calibration gap (πραγματική $-$ προβλεπόμενη συχνότητα): proportional
normalization (αριστερά) έναντι Shin (δεξιά). Το εύρος συρρικνώνεται δραματικά --
οπτική επιβεβαίωση της διόρθωσης του favorite-longshot bias.}
\end{figure}

\textbf{Συμπέρασμα}: το βασικό εύρημα της εργασίας (μείγμα $>$ κάθε μεμονωμένος
bookmaker εκτός του καλύτερου sharp bookmaker· βαθμονόμηση βοηθάει) \textbf{δεν είναι
artifact} της επιλογής de-vig -- επιβεβαιώνεται με δύο ανεξάρτητες μεθόδους.

% =====================================================================
\section{Τρέχουσα κατάσταση \& επόμενα βήματα}
% =====================================================================

Το πλήρες pipeline (download $\to$ επεξεργασία $\to$ πυρήνας αλγορίθμων $\to$
βαθμονόμηση $\to$ 9 επεκτάσεις $\to$ Shin robustness check) τρέχει καθαρά end-to-end. Το
πλήρες κείμενο (θεωρία, υλοποίηση, όλα τα ευρήματα) είναι ήδη σε προχωρημένο στάδιο στο
\file{thesis\_report.tex}. Το ερώτημα που έχει μείνει ανοιχτό είναι \textbf{scope}: πόσα
από τα $\sim$10 follow-up ευρήματα αξίζει να μπουν σε κεντρικό κεφάλαιο έναντι
παραρτήματος, ώστε το κείμενο να έχει καθαρή, εστιασμένη αφήγηση αντί να διασπείρεται.

\end{document}
